\documentclass[]{interact}
\setlength{\voffset}{-.25in}
\sloppy
% packages already loaded by interact.cls: amsmath,amssymb,amsfonts,amsbsy,amsthm,booktabs,epsfig,graphicx,rotating
\usepackage{latexsym}
\usepackage{anyfontsize}
\usepackage{color}
\usepackage{hyperref}
\usepackage{url}
\usepackage{breakurl}
\newcommand{\bburl}[1]{\textcolor{blue}{\url{#1}}}

\makeatletter
\newcommand{\monthyear}[1]{\def\@monthyear{\uppercase{#1}}}
\newcommand{\volnumber}[1]{\def\@volnumber{\uppercase{#1}}}
\makeatother

\begin{document}

\monthyear{[Month Year]}
\volnumber{[Volume, Number]}
\setcounter{page}{1}

\title{A Time-Domain Analogue to Fibonacci Structure via Phase-Gated AR(2) Dynamics:
Reply to Boman on Tissue Fibonacci Patterns and Colonic Crypt Renewal}

\author{
\name{Michael Whiteside\thanks{Email address: mickwh@msn.com}}
\affil{Independent Researcher, Scotland}
}

\maketitle

{\bf Article type}: applications to science

\bigskip

\begin{abstract}
Boman and co-workers have developed an elegant spatial framework in which simple
rules for asymmetric division and niche geometry generate Fibonacci-like
cell-number patterns in renewing tissues [1, 2].
Their tissue-renewal model based on age-structured, asymmetric cell division
produces Fibonacci sequences in time, and their recent five-rule tissue code
for colonic crypts explains how spatial organization is maintained [2].
In this reply, I propose a complementary time-domain analogue: a Phase-Gated
Order-2 Autoregressive (PAR(2)) model for crypt renewal in which Fibonacci-like
structure appears in the dynamics of phase-aligned observables rather than only
in static spatial counts.
I outline the minimal mathematics of the PAR(2) construction, show that
crypt-relevant circadian genes in intestinal organoids are well-described by
stable, complex AR(2) roots, and map the temporal parameters onto Boman's
five spatial rules.
Finally, I derive testable predictions for tuft cell and deep crypt secretory
(DCS) behaviour under circadian disruption, injury, and chronotherapy.
The aim is to offer a simple time-domain model that sits naturally alongside
Boman's spatial code and invites joint experimental tests.
\end{abstract}

\begin{keywords}
AR(2); autoregressive model; circadian; colonic crypt; Fibonacci; phase-gated; tissue renewal; tuft cells
\end{keywords}

%---------------------------------------------------------------------
\section{Context: From Spatial Fibonacci Patterns to a Tissue Code}
%---------------------------------------------------------------------

Fibonacci numbers and related sequences appear widely in biology, from
phyllotaxis to branching and aperiodic order in tissues.
A recent survey by Bormashenko reviews how Fibonacci-like structure can arise
from symmetry, scaling, and growth constraints across multiple systems [3].

Boman et al.\ built on earlier work by Spears and Bicknell-Johnson on
age-structured asymmetric division to propose a tissue-renewal model in which
mature and immature cells follow simple division rules that yield Fibonacci
sequences in time and hierarchical patterns in tissues [1].
In their \emph{Fibonacci Quarterly} paper, they show that asymmetric cell
division with appropriate maturation delays generates Fibonacci cell number
sequences and rational ratios of immature to mature cells over time.

In their recent \emph{Biology of the Cell} article, Boman and colleagues go
further, proposing that a small set of five biological rules---a ``tissue
code''---governs dynamic organization of cells in colonic epithelium [2].
These rules constrain: where stem cells reside; how, when, and in which
direction they divide; how many times progeny may divide; how long
differentiated lineages live; and when and where crypt fission occurs.
Agent-based simulations driven by these rules quantitatively match real crypt
architecture and suggest that the same code may generalize to other tissues [2].

Complementing this, Nguyen, Lausten and Boman provide a comprehensive review of
colonic crypt cellular composition and signalling pathways, including detailed
discussion of LGR5$^{+}$ stem cells, transit-amplifying progenitors,
colonocytes, goblet cells, Paneth-like deep crypt secretory cells (DCS),
enteroendocrine cells, M cells and tuft cells, together with WNT, Notch, BMP
and FGF signalling and their disruption during colorectal cancer and
EMT-driven plasticity [6].
This review offers a biologically rich backdrop for the present time-domain
model.

Taken together, Boman's work therefore offers: (1) a time-based Fibonacci model
for population growth under asymmetric division [1]; (2) a spatial five-rule
code that reproduces static tissue organization [2]; and (3) an integrated
crypt-level description of cell types and signalling, including tuft and DCS
lineages, via Nguyen--Lausten--Boman [6].
This reply asks: what is the simplest time-series model at the crypt level that
could implement and complement this spatial code?

%---------------------------------------------------------------------
\section{Aim of This Reply}
%---------------------------------------------------------------------

The main proposal is a Phase-Gated AR(2) (PAR(2)) model, originally developed
in a more general form for crypt dynamics in the Integrated Temporal-Spatial
Oscillator Hypothesis for Crypt Renewal [11] and applied here specifically to
Boman's Fibonacci tissue code:

\begin{itemize}
  \item The state variable $x_t$ is a phase-aligned crypt-level observable,
        such as the proportion of tuft cells, a weighted renewal index, or a
        composite of several stem and niche markers.
  \item The dynamics obey a second-order autoregression in discrete time:
        $x_t = \varphi_1(\theta_t)\,x_{t-1} + \varphi_2(\theta_t)\,x_{t-2}
        + \varepsilon_t$, where $\theta_t$ is a phase variable (for example,
        circadian or cell-cycle phase), and the coefficients switch or vary
        smoothly with phase.
  \item Under mild constraints on $(\varphi_1, \varphi_2)$, the characteristic
        roots approximate those of a Fibonacci-like recurrence, so that Boman's
        spatial Fibonacci structure is mirrored in temporal return dynamics.
\end{itemize}

The PAR(2) model is deliberately minimal.
It is not intended as a fully mechanistic molecular network, but as a compact
time-domain analogue of Boman's code, with three goals:
(1) to provide a language for stability and overshoot in crypt renewal in
terms of AR(2) roots;
(2) to relate circadian clock disruption and niche signalling to deviations
from Fibonacci-consistent behaviour; and
(3) to generate concrete, falsifiable predictions for tuft cell and DCS
dynamics that can be tested in organoids and animal models.

%---------------------------------------------------------------------
\section{A Phase-Gated AR(2) Model}
%---------------------------------------------------------------------

\subsection{Baseline AR(2) dynamics}

Consider a scalar time series $\{x_t\}$ indexed in discrete steps
(for example, equally spaced circadian sampling times or aligned cell-cycle
checkpoints).
A standard AR(2) process is
$x_t = \varphi_1 x_{t-1} + \varphi_2 x_{t-2} + \varepsilon_t$,
with innovations $\varepsilon_t$ of mean zero and finite variance.
The associated characteristic polynomial
$r^2 - \varphi_1 r - \varphi_2 = 0$
has roots
$r_{1,2} = \bigl(\varphi_1 \pm \sqrt{\varphi_1^2 + 4\varphi_2}\bigr)/2$.
If both roots satisfy $|r_{1,2}| < 1$, the process is stable;
complex-conjugate roots with $|r| < 1$ correspond to a damped discrete-time
oscillator with angular frequency $\arg(r)$ and decay determined by $|r|$.

\subsection{Phase gating}

To introduce phase structure, let $\theta_t \in [0, 2\pi)$ be a phase variable
evolving approximately linearly in time,
$\theta_t = \theta_0 + \Omega t \pmod{2\pi}$,
with $\Omega$ corresponding to a circadian or cell-cycle frequency.
A phase-gated AR(2) model takes
$x_t = \varphi_1(\theta_t)\,x_{t-1} + \varphi_2(\theta_t)\,x_{t-2}
+ \varepsilon_t$,
with $\varphi_1(\cdot)$ and $\varphi_2(\cdot)$ either piecewise constant or
smooth functions of phase.
A minimal two-gate construction uses
$(\varphi_1, \varphi_2) = (\varphi_1^A, \varphi_2^A)$ for $\theta$ in gate A,
and $(\varphi_1^B, \varphi_2^B)$ for $\theta$ in gate B, partitioning the
cycle (for example, subjective day vs.\ night, or G1/S vs.\ G2/M).
Biologically, the gates can encode a proliferative vs.\ quiescent phase, a
high-Wnt vs.\ high-YAP state in crypt bases [8], or a tuft/DCS recruitment
versus maintenance phase conditioned on inflammatory or microbial cues [5, 10].

\subsection{Fibonacci-consistent constraints in time}

Boman's age-structured asymmetric division model yields Fibonacci sequences for
cell counts over discrete time steps by imposing simple rules on division and
maturation [1].
In the scalar AR(2) setting, exact Fibonacci dynamics arise when
$x_n = x_{n-1} + x_{n-2}$, that is, $(\varphi_1, \varphi_2) = (1, 1)$,
whose characteristic roots are the golden ratio $\varphi = (1+\sqrt{5})/2$ and
its conjugate.

In a crypt-level stochastic setting, it is neither necessary nor realistic to
demand exact Fibonacci coefficients.
Instead, I require that the AR(2) coefficients for a phase-aligned observable
$x_t$ lie within a compact, stable region of the parameter plane, that the
associated roots are complex with $|r| < 1$ and phases within a narrow range,
and that under projection into a reduced parameter space, the coefficient cloud
falls near a golden-ratio-consistent manifold associated with Boman-like
asymmetric division rules.
Boman's Fibonacci structure is treated as an attractor in coefficient space
rather than an exact equality.

It is important to note that the Fibonacci point $(\varphi_1,\varphi_2)=(1,1)$
lies \emph{outside} the AR(2) stationarity triangle, with dominant root
$|\lambda| = \varphi \approx 1.618 > 1$.
Consequently, no covariance-stationary gene expression series can operate at
the exact Fibonacci point; all observed gene eigenvalues satisfy $|\lambda|<1$.
The Fibonacci connection is therefore a geometric boundary landmark of the
admissible parameter space, not an attainable operating point.
The stable twin of the Fibonacci characteristic polynomial is $1/\varphi
\approx 0.618$, which lies well inside the stationarity region and is the
biologically relevant reference value.

%---------------------------------------------------------------------
\section{Application to Rhythmic Organoid Data}
%---------------------------------------------------------------------

Stokes et al.\ studied the role of the core clock gene \textit{Bmal1}
(\textit{Arntl}) in intestinal stem cell signalling and tumour initiation [8].
Using organoids derived from mouse intestine, they showed that loss of
\textit{Bmal1} increases self-renewal via YAP1-dependent pathways and disrupts
Wnt signalling; transcriptomic profiling was deposited as GEO series
GSE157357 [9].

\textbf{Limitation: model--data mismatch.}
Boman's five-rule tissue code was developed for \emph{human large intestinal}
crypts [2], while Stokes et al.\ used \emph{mouse small intestinal}
organoids [8].
Both the species difference (human vs.\ mouse) and the anatomical region
difference (large vs.\ small intestine) are genuine limitations.
Colon and small intestine differ in LGR5$^{+}$ stem cell density, villus
structure, crypt length, and cancer susceptibility; and mouse and human crypt
architecture, while broadly conserved, are not identical.
The AR(2) fits below should therefore be read as a proof-of-concept
demonstration of the method on the best available rhythmic organoid dataset,
not as a validated quantitative test of Boman's human-colon model.
Direct validation against a human large-intestinal, circadian time-series
dataset remains an important open objective.

For an illustrative PAR(2) analysis, consider phase-ordered expression of four
crypt-relevant genes from GSE157357: \textit{Lgr5} (stem cell marker),
\textit{Arntl} (Bmal1; core clock), \textit{Per2} (clock output), and
\textit{Axin2} (Wnt pathway readout).
Using a simple least-squares AR(2) fit to $z$-scored log-expression values
ordered by circadian phase, I obtain the characteristic roots shown in
Table~\ref{tab:ar2}.

\begin{table}[h]
\tbl{Illustrative AR(2) fits to selected genes in GSE157357 [8, 9].
Complex-conjugate roots with modulus $|r|<1$ indicate stable, damped
oscillations in the phase-ordered signal.
Root magnitudes well below 1 correspond to strongly damped behaviour.}
{\begin{tabular}{lccc}
\toprule
\textbf{Gene} & \textbf{Roots} & \textbf{Max $|r|$} & \textbf{Stability} \\
\midrule
\textit{Lgr5}  & $0.049 \pm 0.332j$ & 0.336 & Stable \\
\textit{Arntl} & $0.017 \pm 0.061j$ & 0.063 & Stable \\
\textit{Per2}  & $-0.092 \pm 0.513j$ & 0.521 & Stable \\
\textit{Axin2} & $-0.008 \pm 0.464j$ & 0.464 & Stable \\
\bottomrule
\end{tabular}}
\label{tab:ar2}
\end{table}

All four examples yield complex-conjugate roots with $|r| \ll 1$, indicating
that the phase-ordered expression behaves as a damped oscillator around a mean
level.
Within the PAR(2) framework, these gene-level AR(2) fits are interpreted as
projections of a lower-dimensional crypt-level oscillator into different
observables.
A natural next step would be to move from univariate AR(2) to multivariate
VAR(2) models for panels of genes, fit phase-dependent coefficients, and test
whether renewal-relevant composites cluster near a Fibonacci-consistent manifold
under wild-type conditions and deviate under Bmal1 loss or environmental
circadian disruption.

%---------------------------------------------------------------------
\section{Tuft Cells, Deep Crypt Secretory Cells, and Long-Lived Lineages}
%---------------------------------------------------------------------

The PAR(2) framework requires three specific inputs to be testable: (i) a
rhythmic time-series dataset from crypt-relevant tissue, (ii) a long-lived
cell type whose dynamics accumulate temporal perturbations over many cycles,
and (iii) a niche-supporting cell type whose sustained signalling activity
could implement the memory kernels $\alpha_1$ and $\alpha_2$.
The GSE157357 organoid dataset [8, 9] is chosen because it is the only
publicly available circadian time-series profiling crypt stem cell behaviour
across wild-type and clock-disrupted conditions simultaneously, making it
uniquely suited to test the PAR(2) kernel architecture.
Tuft cells are chosen because they are among the longest-lived epithelial
cells in the crypt (lifespan $\geq 28$ days vs.\ the 3--5 day average
turnover), meaning they accumulate the effects of temporal perturbations
across many renewal cycles and are therefore the most sensitive available
readout of sustained eigenvalue shifts [10].
DCS cells are chosen because their anatomical position at the crypt
base---intermingled with LGR5$^{+}$ stem cells and providing continuous WNT,
EGF and Notch ligand supply---makes them the natural physical implementors of
the sustained memory kernels $\alpha_1$ (fast, NOTCH-mediated) and $\alpha_2$
(slow, WNT/AXIN2-mediated) in the PAR(2) model [6, 7].
Long-lived lineages more generally are those most likely to accumulate
perturbations in AR(2) eigenvalues over many cell cycles and therefore provide
the most sensitive readout of sustained temporal disruption.

Tuft cells are rare chemosensory epithelial cells with important roles in
immunity and disease.
Westphalen et al.\ used lineage tracing to show that a DCLK1$^{+}$ tuft cell
population in the intestine is long-lived, contributes to regeneration after
injury, and can serve as a colon cancer-initiating cell following APC mutation
and inflammatory insult [10].
Subsequent work has expanded their role in antimicrobial and anti-parasitic
defence and in shaping mucosal immunity [5].

Within the colonic crypt, Nguyen, Lausten and Boman highlight tuft cells,
M cells and enteroendocrine cells as important contributors to tumour
heterogeneity and treatment resistance, particularly via
epithelial--mesenchymal transition and cellular plasticity [6].
Tuft cell expansion is noted as one component of tumour--immune crosstalk and
microenvironmental disruption in colorectal cancer [6].

DCS cells have been characterised as Reg4$^{+}$ niche-supporting cells
intermingled with LGR5$^{+}$ stem cells at the crypt base [6, 7].
Circadian control of stem cell niche components, including regulation of
\textit{Lgr5} expression by BMAL1 and RNA-binding protein MEX3A [4], suggests
that DCS cell function may be phase-modulated, consistent with a role as
implementors of the PAR(2) slow memory kernel $\alpha_2$.

From a tissue-code perspective, long-lived tuft cells and related DCS lineages
sit at the intersection of division limits and lifespan (rules on how many
times and how long cells persist), niche signals (for example, IL-25, microbial
metabolites, Hippo/Wnt balance), and injury responses [5, 6, 10].

Under perturbations known to affect tuft/DCS biology---for example, helminth
infection, chronic inflammation, microbial changes, or deliberate timing
interventions [5, 6, 10]---the PAR(2) model predicts systematic shifts in
$(\varphi_1, \varphi_2)$ and therefore in root magnitudes and phases.

%---------------------------------------------------------------------
\section{Mapping Boman's Five Rules onto PAR(2) Dynamics}
%---------------------------------------------------------------------

Boman's five rules for colonic crypt organization can be mapped onto PAR(2)
parameters at a coarse level.
Rule~1 (spatial positioning of stem cells) constrains the effective pool size
and interaction network feeding into $x_t$, narrowing the allowable region of
$(\varphi_1, \varphi_2)$ and preventing runaway growth.
Rule~2 (order and direction of divisions) specifies directed movement and
differentiation of daughters along the crypt axis, appearing in
phase-dependent coefficients $\varphi_1(\theta)$, $\varphi_2(\theta)$ that
encode preferred advance directions.
Rule~3 (division limits) places an upper bound on contributions from older
generations, pushing the system toward low-order autoregressions with
$|r| < 1$.
Rule~4 (lifespan of differentiated cells) sets how quickly perturbations
decay: shorter lifespans map to smaller $|r|$; extended lifespans push $|r|$
toward 1.
Rule~5 (crypt fission) corresponds to a branching event, which can be modelled
as a change-point where coefficients switch regime or as a stochastic
bifurcation when $|r|$ or innovation variance crosses a critical
threshold [2].

In this view, the five rules act as constraints that restrict which AR(2)
parameter sets are allowed.
Fibonacci structures identified by Boman emerge as special cases where the
resulting dynamics stay near golden-ratio-consistent manifolds, both in space
and time [1, 2].

A critical clarification: the Fibonacci recurrence $x_n = x_{n-1}+x_{n-2}$
corresponds to $(\varphi_1,\varphi_2)=(1,1)$, which lies \emph{outside} the
AR(2) stationarity triangle with dominant root $|\lambda|\approx 1.618>1$.
This means that exact Fibonacci dynamics would be explosive, not stable.
Living tissues operate in the heavily damped interior of the stationarity
triangle: the observed root magnitudes in Table~\ref{tab:ar2} range from
$|r|=0.063$ (\textit{Arntl}) to $|r|=0.521$ (\textit{Per2}), far from the
explosive boundary.
The biological content of the Fibonacci analogy is therefore not that tissue
renewal grows like the Fibonacci sequence---it does not---but that homeostatic
regulation holds the AR(2) coefficients near the stable twin of the Fibonacci
eigenvalue, $1/\varphi\approx 0.618$, which is well inside the admissible
region.
The five rules of Boman's tissue code are precisely the biological mechanisms
that enforce this heavily damped, sub-unit-root regime.

To make the Fibonacci-consistent manifold concrete: it corresponds to the
subset of the stable AR(2) stationarity triangle---the biologically admissible
parameter space---in which the coefficient pair $(\varphi_1, \varphi_2)$ lies
close to the Fibonacci boundary landmark at $(1,1)$.
In the $(\varphi_1, \varphi_2)$ plane, the Fibonacci point $(1,1)$ sits at the
boundary of the stable region and serves as a geometric reference: pairs whose
dominant root $|r|$ falls in a narrow band below unity, consistent with the
observed damped-oscillator behaviour in Table~\ref{tab:ar2}, define the
admissible neighbourhood.
The exact extent of this neighbourhood is determined by the division limit
(Rule~3) and lifespan (Rule~4) constraints of Boman's tissue code, which
together bound $|r|$ from above and below.
A formal characterisation and visualisation of this region is provided in the
companion platform at \bburl{https://par2-discovery-engine.replit.app}
(Root-Space Geometry page).

%---------------------------------------------------------------------
\section{Testable Predictions and Experimental Designs}
%---------------------------------------------------------------------

The PAR(2) framework is simple enough to support explicit, falsifiable
predictions.
For example, circadian disruption (for example, Bmal1 loss or environmental
misalignment) should increase variance in $(\varphi_1, \varphi_2)$ across
crypts and broaden the distribution of root magnitudes $|r|$ and phases
$\arg(r)$, shifting the coefficient cloud away from a Fibonacci-consistent
region [8].
Long-lived DCLK1$^{+}$ tuft cells should exhibit a damped AR(2)-like overshoot
after injury, with longer return times (larger $|r|$) in chronic inflammatory
contexts [10].
Chronotherapy that restores phase alignment should narrow dispersion in
$(\varphi_1, \varphi_2)$, shrink the spread of $|r|$, and decrease
inter-crypt variance in tuft/DCS observables.
Finally, agent-based simulations implementing Boman's five rules could be used
to generate synthetic time-series of crypt observables and verify that stable
Fibonacci-like regimes produce complex roots with $|r| < 1$ and coherent
phase, while perturbations to individual rules map to distinct deformations in
AR(2) coefficient space [2].

%---------------------------------------------------------------------
\section{Discussion and Outlook}
%---------------------------------------------------------------------

Boman's contributions to Fibonacci biology and tissue organization provide a
powerful framework for understanding how simple rules can generate complex,
ordered patterns in tissues [1, 2].
The Phase-Gated AR(2) model proposed here is intentionally modest: a
low-dimensional time-domain representation whose parameters can be estimated
from existing and future data, and which makes explicit contact with Fibonacci
recurrences via its characteristic roots.
The present work therefore sits at the intersection of Boman's tissue-renewal
matrices [1, 2], the crypt-wide cellular and signalling landscape described by
Nguyen, Lausten and Boman [6], and the Integrated Temporal-Spatial Oscillator
Hypothesis for Crypt Renewal [11], specialising the latter to the specific case
of colonic crypt dynamics and Fibonacci tissue codes.

In Boman's Fibonacci Quarterly model, asymmetric division and maturation rules
are encoded in a low-dimensional transition matrix whose eigenvalues generate
Fibonacci sequences in time and stable cell-type ratios [1].
In the present work, the proposed Phase-Gated AR(2)/VAR(2) framework plays an
analogous role at the level of crypt observables: the AR/VAR coefficients form
a companion matrix whose eigenvalues encode growth, decay and oscillation of
renewal-related variables.
Rather than demanding exact Fibonacci recurrences, Boman's solution is treated
as a reference point in this eigenvalue space, and the question is whether
homeostatic crypts lie near a ``Fibonacci-consistent'' region, with circadian
disruption, injury or oncogenic mutation manifesting as systematic excursions
away from that region.

For a journal like \emph{The Fibonacci Quarterly}, the main message is that
Fibonacci-like structure in tissues need not be confined to static spatial
counts.
It can also appear in the temporal structure of phase-aligned return dynamics,
with golden-ratio-related constraints shaping the AR(2) coefficient space.
Boman's five-rule tissue code can be viewed as defining a family of admissible
PAR(2) models, with pathological states corresponding to excursions away from
this family.
The Fibonacci-consistent manifold has been defined above (Section~6) as the
subset of the stable AR(2) stationarity triangle---the biologically admissible
parameter space---in which the coefficient pair lies close to the Fibonacci
boundary landmark at $(1,1)$, with the admissible neighbourhood bounded by the
division-limit and lifespan constraints of Rules~3 and~4.
Future work should apply multivariate PAR(2) and VAR(2) models to circadian and
injury time-series from crypts and organoids, and integrate spatial Boman-style
simulations with time-domain summaries to build a unified spatio-temporal
Fibonacci framework in which the manifold boundary can be validated against
experimental perturbations.

%---------------------------------------------------------------------
\section{Data Availability}
%---------------------------------------------------------------------

The Phase-Gated AR(2) / temporal-spatial oscillator framework is archived on
Zenodo under DOI 10.5281/zenodo.17366927 [11].
The organoid transcriptomic data re-analysed here are publicly available from
the NCBI Gene Expression Omnibus under accession GSE157357 [9], as originally
reported in Stokes et al.\ [8].
No additional primary datasets were generated for this reply.

%---------------------------------------------------------------------
\section{AI Use Declaration}
%---------------------------------------------------------------------

In accordance with the Taylor \& Francis AI Policy, the author declares the
following.
Generative AI tools were used during the preparation of this manuscript.
Specifically, ChatGPT (model: GPT-5.1 Thinking, OpenAI, 2025) was employed to
assist with: (i) editing and refining the English prose for clarity, concision
and consistency; (ii) \LaTeX{} and document structuring, equation typesetting,
and consistency checks; (iii) support in formulating and refining the
PAR(2)/VAR(2) mathematical framework, including notation, companion-matrix
forms and explanatory text; and (iv) guidance on exploratory analysis of
publicly available transcriptomic data (for example, structuring time-series
inputs, suggesting AR(2) fitting strategies and summarising root behaviour)
based on datasets and results specified by the author.
These tools were used as aids to drafting, calculation and exposition rather
than as independent scientific authors.
All scientific hypotheses, modelling choices, and formal conclusions presented
in this article were reviewed, curated and integrated by the human author, who
has critically assessed all AI-assisted content and takes full responsibility
for the accuracy, integrity and originality of the work.

%---------------------------------------------------------------------
% Disclosure statement — required by journal (added per editorial instruction)
%---------------------------------------------------------------------
\bigskip
\noindent\textbf{Disclosure statement.}
The author filed a UK patent application (GB2518973.9) related to the
AR(2) eigenvalue methodology described in this work on 13 November 2025,
three days prior to the original manuscript submission on 16 November 2025.
The author declares no other conflicts of interest.

%---------------------------------------------------------------------
\begin{thebibliography}{11}
%---------------------------------------------------------------------
% References in alphabetical order by last name, Fibonacci Quarterly style.
% Format: Author(s). Title. J. Abbrev. Volume(Issue):pages, Year. doi:...

\bibitem{boman2017fibonacci}
Boman, B.M., Dinh, T.Q., Decker, K., Emerick, B., Raymond, C., and
Schleiniger, G.
Why do Fibonacci numbers appear in patterns of growth in nature? A model for
tissue renewal based on asymmetric cell division.
Fibonacci Quart. 55(5):30--41, 2017.

\bibitem{boman2025dynamic}
Boman, B.M. et al.
Dynamic organization of cells in colonic epithelium is encoded by five
biological rules.
Biol. Cell. 117(7):e70017, 2025.
doi:10.1111/boc.70017.

\bibitem{bormashenko2022fibonacci}
Bormashenko, E.
Fibonacci sequences, symmetry and order in biological patterns, their sources,
information origin and the Landauer principle.
Biophysica. 2(3):Art.\ 27, 2022.
doi:10.3390/biophysica2030027.

\bibitem{cheng2023bmal1}
Cheng, L.-T. et al.
Core clock gene BMAL1 and RNA-binding protein MEX3A collaboratively regulate
Lgr5 expression in intestinal crypt cells.
Sci. Rep. 13:17597, 2023.
doi:10.1038/s41598-023-44997-5.

\bibitem{hendel2022tuft}
Hendel, S.K. et al.
Tuft cells and their role in intestinal diseases.
Front. Immunol. 13:822867, 2022.
doi:10.3389/fimmu.2022.822867.

\bibitem{nguyen2025colonic}
Nguyen, A.L., Lausten, M.A., and Boman, B.M.
The colonic crypt: cellular dynamics and signaling pathways in homeostasis and
cancer.
Cells. 14(18):1428, 2025.
doi:10.3390/cells14181428.

\bibitem{schumacher2023dcs}
Schumacher, M.A.
The emerging roles of deep crypt secretory cells in colonic physiology.
Am. J. Physiol. Gastrointest. Liver Physiol. 325(6):G493--G500, 2023.
doi:10.1152/ajpgi.00093.2023.

\bibitem{stokes2021bmal1}
Stokes, K. et al.
The circadian clock gene, Bmal1, regulates intestinal stem cell signaling and
represses tumor initiation.
Cell. Mol. Gastroenterol. Hepatol. 12(5):1847--1872, 2021.
doi:10.1016/j.jcmgh.2021.08.001.

\bibitem{stokes2021geo}
Stokes, K. et al.
Transcript profiling dataset: The circadian clock gene, Bmal1, regulates
intestinal stem cell signaling and represses tumor initiation.
GEO accession GSE157357, NCBI, 2021.

\bibitem{westphalen2014tuft}
Westphalen, C.B. et al.
Long-lived intestinal tuft cells serve as colon cancer-initiating cells.
J. Clin. Invest. 124(3):1283--1295, 2014.
doi:10.1172/JCI73434.

\bibitem{whiteside2025integrated}
Whiteside, M.
Integrated temporal-spatial oscillator hypothesis for crypt renewal.
Zenodo preprint, 2025.
doi:10.5281/zenodo.17366927.

\end{thebibliography}

\end{document}
